\newcommand{\latticeToMLB}{%
\begin{figure}[t]
\centering
\resizebox{0.85\textwidth}{!}{%
\begin{tikzpicture}[
  >=Latex,
  line cap=round, line join=round,
  thick,
]
  % --- styles
  \colorlet{nodegreen}{green!60!black}
  \tikzset{
    gedge/.style={draw=nodegreen, line width=1.8pt},
    gnode/.style={draw=nodegreen, fill=white, line width=1.8pt},
    arrow/.style={->, line width=1.8pt, draw=black},
    lbl/.style={font=\Large},
  }

  % --- coordinates (tree on the left)
  \coordinate (n1) at (0,2.8);
  \coordinate (n2) at (-1.6,1.4);
  \coordinate (n3) at ( 1.6,1.4);
  \coordinate (n4) at ( 0.8,0.0);
  \coordinate (n5) at ( 2.4,0.0);

  % --- green edges
  \draw[gedge] (n1) -- (n2);
  \draw[gedge] (n1) -- (n3);
  \draw[gedge] (n3) -- (n4);
  \draw[gedge] (n3) -- (n5);

  % --- green open nodes
  \draw[gnode] (n1) circle [radius=0.30];
  \draw[gnode] (n2) circle [radius=0.28];
  \draw[gnode] (n3) circle [radius=0.28];
  \draw[gnode] (n4) circle [radius=0.26];
  \draw[gnode] (n5) circle [radius=0.26];

  % --- labels (node numbers)
  \node[lbl, anchor=south] at ($(n1)+(0,-0.85)$) {$p$};
  \node[lbl, anchor=north] at ($(n2)+(0,-0.25)$) {$x$};
  \node[lbl, anchor=north] at ($(n3)+(0,-0.25)$) {$q$};
  \node[lbl, anchor=north] at ($(n4)+(0,-0.25)$) {$y$};
  \node[lbl, anchor=north] at ($(n5)+(0,-0.25)$) {$z$};

  % --- middle arrow + label
  \coordinate (Astart) at (3.6,1.4);
  \coordinate (Aend)   at (6.6,1.4);
  \draw[arrow] (Astart) -- (Aend);
  \node[font=\bfseries\Large] at ($(Astart)!0.5!(Aend)+(0,0.6)$) {\ac{mlb}};

  % --- right-side explanation (siblings + uncle)
  \node[anchor=west, font=\Large] at (7.1,1.8) {$(y,z,x)$};
  \node[anchor=west, font=\Large] at (7.1,1.1)
    {\small not $(y,z,q)$ because $q$ is the parent/meet of $y$ and $z$;};
  \node[anchor=west, font=\Large] at (7.1,0.5)
    {\small not $(x,q,p)$ because $p$ is the parent/meet of $x$ and $q$.};

\end{tikzpicture}%
}
\caption{Building \ac{mlb} constraints using the ``siblings + uncle'' rule, where siblings $y,z$ are connected to the same parent node $q$ and $x$ is the uncle node.}
\label{fig:concept_lattice2mlb}
\end{figure}%
}

\subsection{\acl{mlb} Constraints}
\label{sec:foundations_mlb}

Clustering structures unlabelled data, which is particularly useful in real-world applications where labeled data is scarce and annotation is costly. 
Clustering algorithms group data points according to a similarity criterion, typically minimizing intra-cluster distances while maximizing inter-cluster distances.
Prior knowledge about the dataset can be incorporated through constraints. 
The simplest forms are pairwise \textit{must-link} and \textit{cannot-link} constraints, enforcing that two instances must or must not belong to the same cluster. 
However, pairwise constraints generally apply only at one level of granularity.

Hierarchical clustering instead organizes data into nested groups at different levels of granularity~\cite{bade_hierarchical_2014}. 
To incorporate prior knowledge in this setting, \ac{mlb} constraints specify the relative agglomeration order of three data points $d_x$, $d_y$, $d_z$. 
An \ac{mlb} constraint
\begin{eqnarray}
    MLB_{x y z} = (d_x, d_y, d_z)
    \label{eq:mlb_xyz}
\end{eqnarray}
states that $d_x$ and $d_y$ should be merged before either is merged with $d_z$, where $d_x$ and $d_y$ are symmetric. 
Let $H = (C, \prec, \text{root})$ denote a cluster hierarchy, where $C$ is a set of clusters, $\prec$ is a strict partial order over $C$ representing the parent-child relation, and $\text{root} \in C$ is the root cluster~\cite{bade_hierarchical_2014}. 
Definitions~\eqref{def:mlb_implication} and \eqref{def:mlb_counterexample} formalize the idea that $d_x$ and $d_y$ are more closely related in the hierarchy than $d_x$ and $d_z$. 
\autoref{fig:hierarchy_to_mlb} illustrates valid \ac{mlb} constraints for a toy example.

\begin{definition}[\ac{mlb} Implication]
    \label{def:mlb_implication}\\
    Given \ac{mlb} constraint $MLB_{x y z}$:
    $\forall c \in H: d_x \in c \wedge d_z \in c \rightarrow d_y \in c$
\end{definition}
\begin{definition}[\ac{mlb} Counter Example]
    \label{def:mlb_counterexample}\\
    Given \ac{mlb} constraint $MLB_{x y z}$:
    $\exists c \in H: d_x \in c \wedge d_y \in c \wedge d_z \notin c$
\end{definition}

\hierarchyToMLB{}

% A \ac{mlb} constraint for hierarchical clustering can be transformed into multiple pairwise constraints for partitional clustering by adding one must-link constraint between $d_x$ and $d_y$ and two cannot-links between $d_x, d_z$ and $d_y, d_z$~\cite{bade_hierarchical_2014}.

% \subsubsection{Concept hierarchies} 
% Concept hierarchies structure information into hierarchies and facilitate working with knowledge-bases~\cite{cimiano_learning_2005}.
